\clearpage
\section{Literature review}
\label{sec:literature}

Multi-omics integration sits at the intersection of computational biology and representation learning. This chapter reviews classical statistical approaches, deep multi-omics models, the shared--private contrastive line of work that includes MOCSS and OMIDIENT, and broader machine-learning critiques of claimed disentanglement. The chapter closes by stating the gap addressed in this thesis.

\subsection{Classical integrative clustering and factorisation}

Early multi-omics subtype discovery largely relied on linear latent-variable or graph-fusion models. iCluster~\citep{shen2009integrative} poses integrative clustering as a joint Gaussian latent-variable model: a common discrete cluster assignment generates modality-specific continuous loadings, so that subtype structure is forced to be shared while residual variation remains view-local. The method is statistically clear and remains a common baseline, but scalability and nonlinear interactions are limited.

MOFA~\citep{argelaguet2018multi} and related factor models extend the idea to continuous latent factors with sparsity and likelihood choices adapted to each omics type. Factors can be shared or private depending on which views they load on, providing an interpretable decomposition without requiring a single clustering objective. JIVE~\citep{lock2013jive} makes the shared--individual partition explicit at the matrix level by decomposing a multi-block data set into joint structure, individual structure, and residual noise. These methods establish the conceptual vocabulary used throughout this thesis: shared structure should explain cross-view agreement, while private structure should explain leftover within-view variance.

Graph-based approaches take a complementary route. Similarity Network Fusion~\citep{wang2014snf} constructs a patient similarity network per modality and iteratively fuses them so that local neighbourhoods that agree across views are reinforced. Spectral clustering on the fused graph then yields subtypes~\citep{ng2002spectral,vonluxburg2007tutorial}. SNF does not learn parametric encoders, but it demonstrates that cross-view agreement is naturally expressed as geometry on samples rather than as independent per-view reconstructions. That geometric intuition reappears in Method~3 of this thesis.

Benchmarking studies have emphasised that no single classical method dominates all cancers and metrics~\citep{rappoport2021evaluation,cantini2021benchmarking}. Performance depends on preprocessing, the number of clusters, and whether the biological ground truth is itself contested. Consequently, modern work increasingly reports multiple metrics and stress-tests robustness across seeds and hyperparameter choices~\citep{baiao2025technical}.

A further lesson from the classical literature is that ``shared'' has always been an operational concept. In iCluster it is the discrete assignment; in MOFA and JIVE it is the set of factors or matrix components with multi-view support; in SNF it is the fused neighbourhood relation. None of these definitions reduces to ``whatever remains after training an autoencoder with an orthogonality penalty.'' That historical precision motivates the probe-centred evaluation used later in this thesis.

\subsection{Deep multi-omics representation learning}

Deep models replace linear factors with nonlinear encoders, typically autoencoders or variational autoencoders~\citep{kingma2013auto,Goodfellow-et-al-2016}. The appeal is capacity: nonlinear maps can capture interactions among genes, CpG sites, and microRNAs that linear factor models miss. Surveys of multi-omics machine learning highlight both the promise and the risk of opaque embeddings that are hard to validate biologically~\citep{hasin2017multi,xiao2022multi,jiang2022big,mohr2024navigating}. Large-scale cancer genomics programmes have accelerated the availability of matched multi-omics, but they have also increased the temptation to treat any embedding that clusters by clinical labels as scientifically resolved~\citep{jiang2022big}.

Architecturally, deep integration strategies fall into early, intermediate, and late fusion. Intermediate fusion---learning per-view encoders and then combining latent codes---is the dominant pattern for unsupervised subtyping because it respects modality heterogeneity while still enabling joint clustering. Early concatenation ignores scale and noise differences; late fusion of separately clustered views can miss coordinated but weak cross-omics signal. Intermediate shared--private models attempt to keep the best of both worlds by allocating capacity to cross-view and view-local factors inside a common training loop~\citep{xiao2022multi}.

Contrastive learning has become a popular alignment tool in this setting: positive pairs are matched samples across views, and InfoNCE-style objectives pull their embeddings together~\citep{chen2020simple}. When applied to multi-omics, contrastive losses encourage cross-view consistency of the aligned subspace, but they do not by themselves define what should remain private. If the shared encoder can still reconstruct local assay detail, contrastive alignment may merely synchronise redundant copies of mixed information.

Several deep shared--private designs have been proposed. The common template encodes each view into~$S_m$ and~$P_m$, applies a cross-view consistency loss on~$S$, and reconstructs~$X_m$ from some combination of~$S_m$ and~$P_m$. The details of reconstruction and independence regularisation determine whether the intended semantics hold. Weak orthogonality between~$S_m$ and~$P_m$ is often used as a cheap surrogate for disentanglement, yet elementwise product penalties do not imply subspace orthogonality or statistical independence~\citep{zbontar2021barlow,bardes2022vicreg}. Principal-component diagnostics of omics matrices already show that dominant variance directions need not coincide with biologically annotated structure~\citep{yeung2001principal}; likewise, dominant reconstructive directions in a private code need not be ``private'' in the scientific sense.

\subsection{MOCSS and OMIDIENT}

MOCSS~\citep{chen2023mocss} is a representative deep shared--specific model for cancer subtyping. Each modality is mapped by an autoencoder into shared and specific latents. Two design choices are especially consequential. First, both the shared decoder and the specific decoder reconstruct the full input~$X_m$, so each code is pressured to be individually sufficient for reconstruction. Second, shared codes are projected by a multilayer perceptron and aligned with a contrastive (InfoNCE) loss across modalities, while a weak elementwise orthogonality term discourages co-activation of~$S_m$ and~$P_m$. Empirically, MOCSS reports strong unsupervised subtype recovery on TCGA cohorts. The present thesis revisits that success through role-aware probes and finds that high concatenation NMI can coexist with private codes that retain---and sometimes dominate---subtype structure. Equally telling is the cross-view reconstruction pattern: predicting a modality~$X_A$ from the shared codes of the other modalities yields approximately the same Ridge~$R^{2}$ as predicting~$X_A$ from those modalities' private codes, so the cross-view gap~$\Delta\mathrm{X}\,R^{2}$ stays near zero. If~$S$ and~$P$ truly played distinct roles, foreign shared codes should systematically outperform foreign private codes at explaining another view; the near-parity of the two predictors is further evidence that the named shared and private latents are not what the architecture labels suggest.

From a modelling standpoint, dual full reconstruction is in tension with the residual interpretation of private codes. If~$P_m$ must reconstruct~$X_m$ alone, it has little incentive to leave cross-view signal exclusively to~$S_m$. Orthogonality that penalises only per-dimension co-activation after normalisation cannot close that loophole: nonlinear mixing and subspace overlap remain possible. The contrastive term improves cross-view agreement of projected~$S$, but agreement is compatible with both pure shared biology and synchronised leakage of private detail.

OMIDIENT~\citep{Safinianaini_Välimäki_Bresson_Gorbonos_Rajamäki_Aaltonen_Marttinen_2025} retains a MOCSS-style shared pathway but replaces the private branch with a product-Dirichlet variational encoder--decoder. Its primary design goal is to improve the \emph{private/specific} representation---exploiting compositional Dirichlet structure and modelling view-specific uncertainty more faithfully than a deterministic autoencoder---and the original work reports gains on clustering, classification, and missing-data reconstruction that demonstrate the value of that private-branch change. The paper's central claim is therefore about a better private model within a shared--specific template, not about establishing a new operational definition of sharedness or about proving that~$S$ and~$P$ are information-theoretically distinct. In principle, a well-behaved private Dirichlet code should absorb residual assay-specific variation while leaving cross-view subtype biology in~$S$. When the same BRCA checkpoints are re-examined with the role probes of this thesis, multi-seed evaluation reveals near-zero $\mathrm{concat}(P)$ NMI together with strongly negative own-view $X\leftarrow P$ Ridge~$R^{2}$. That pattern is consistent with private-path collapse under the probe suite used here, rather than with a verified shared--private theory win. The point is not that OMIDIENT lacks scientific value---on the contrary, it advances private modelling and remains a strong integration baseline---but that low private-subtype NMI alone should not be read as proof that the shared/private labels have been validated. Moreover, single-seed subtype scores that look favourable can shrink under multi-seed aggregation, underscoring the need for seed-robust reporting~\citep{baiao2025technical}. Because OMIDIENT's shared branch is still MOCSS-like, improvements confined to the private likelihood do not automatically repair the definition of sharedness; that complementary question is the focus of the present thesis.

\subsection{Shared--private learning beyond multi-omics}

Outside genomics, representation learning has long distinguished factors that should be invariant across views from factors that should remain style- or domain-specific~\citep{bengio2013representation,Goodfellow-et-al-2016}. Contrastive methods such as SimCLR~\citep{chen2020simple} maximise agreement under augmentation, while non-contrastive methods such as Barlow Twins~\citep{zbontar2021barlow} and VICReg~\citep{bardes2022vicreg} control redundancy through covariance and variance constraints. These ideas are transferable to multi-omics if one treats modalities as views: shared codes should agree across views, and private codes should explain residuals without duplicating the shared subspace.

Critiques of claimed disentanglement emphasise that independence penalties, reconstruction losses, and clustering metrics measure different properties~\citep{martens2024disentangling}. A model can look disentangled under one probe and entangled under another. In multi-omics, an especially misleading failure mode is to treat low private-subtype NMI as proof of successful separation: private codes may be collapsed, or they may encode complementary biology that is subtype-related yet not redundant with~$S$. Levine et al.~\citep{levine2013integrated} and related clinical multi-omics studies illustrate that assay-specific correlates of subtype are scientifically meaningful; discarding them by force can harm both interpretation and integration. Navigating multi-omics method choice therefore requires matching the inductive bias to the scientific question, not only to a leaderboard metric~\citep{mohr2024navigating}.

\subsection{Evaluation practice and the research gap}

Technical surveys caution that preprocessing, feature selection, and metric choice can dominate apparent method rankings~\citep{baiao2025technical,rappoport2021evaluation,cantini2021benchmarking}. Standard unsupervised scores---NMI, ARI, clustering accuracy after Hungarian matching~\citep{Kuhn1955Hungarian}, and neighbourhood accuracy---remain necessary for comparing subtype recovery, but they are silent on latent semantics. For shared--private models the evaluation gap is sharper still: standard leaderboards report clustering on $\mathrm{concat}(P,S)$ or on~$S$ alone, but rarely quantify whether foreign private codes predict another modality better than foreign shared codes. Without such probes, one cannot tell whether ``shared'' is shared.

The gap addressed by this thesis is therefore twofold. Methodologically, existing deep shared--private pipelines do not sufficiently constrain~$S$ and~$P$ to their intended roles when both codes are allowed to reconstruct~$X_m$ under weak independence regularisation. Evaluatively, high unsupervised subtype NMI has been over-interpreted as evidence of successful decomposition. The following chapters develop Methods~1--3---predictability constraints, joint reconstruction with VICReg, and spectral shared targets---and a probe-centred protocol to close that gap on BRCA, with KIRC reserved for a later revision.
